% Fresh proof, 2026-10-22; all frozen reconstruction inputs are unchanged.
\section{R3 and R4 at infinitesimal \(e\), with all source ratios and tests}
\label{sec:new-r34-all-tests}

\paragraph{Statement and exact domain.}
Use the original common-II stationary source, its six native weights
$n,m,C,A,D,\alpha$, all its strict source gaps, $d=1$,
$\lambda=\beta=0$, and $n\geq0$. Suppose $B_s,a_sB_s,t_sB_s$
are positive infinitesimals, $t_s,h$ are positive bounded units,
and $e=1-r$ is positive infinitesimal. In addition impose either
\[
 \begin{array}{ll}
 \mathrm{R3}:&1-h\text{ is infinitesimal},\\
 \mathrm{R4}:&1-h\text{ is a positive unit and }
       1+a_sB_sr-(1-h)/t_s\text{ is infinitesimal}.
 \end{array}                                                    \tag{EN1}
\]
There is no condition on the size of $a_s$, $a_s^2B_s$,
$a_s e$, or their ratios. The conclusion is finite separated tensor
SOS for the entire original parent at every independent projective
testing pair and every state, including signs, axes, and literal $n=0$.
The already recovered positive-unit-$e$ proofs then complete the
whole original R3 and R4 rows. No other weak-C2 row is inferred.

The source and testing fields are independent real closed extensions
$K_c,K_t\supset\mathbb R$. A bounded coefficient means an element
of $B_c\otimes B_t$; its residue is the homomorphism to $\mathbb R$.
The strict bounded-Gram lemma is used only for matrices with positive
definite real residue. Row entries may be arbitrarily large.
Every inverse and radical below is pure source, pure testing, or a
product of such factors.

\subsection{One exact compact source for both ends}

Use the exact compact variables of the reconstructed actual-source
dictionary RSR2--RSR6:
\[
\begin{aligned}
c_s&=1+a_sB_sr,\quad R_s=1+ha_sB_s,\quad
j=a_sc_s/R_s,\quad B=B_sR_s/c_s,\quad T=t_sc_s/R_s,\\
p&=a_s e/R_s,\quad H=(1-h)/(t_sc_s),\qquad
x=jB=a_sB_s,\quad y=TB=t_sB_s,\quad q_z=jT .
\end{aligned}                                                    \tag{EN2}
\]
Thus $0<p<\min(1,j)$, $0<H<1$, and
\[
B,x,y\in\mathrm I_c,\quad T\in\mathrm U_c,\quad
1-HT\in\mathrm U_c,\quad p/j\in\mathrm I_c.              \tag{EN3}
\]
The last assertion is exact source entry:
$e=p(1+x)/(j+px)\sim p/j$.
In R3, $H$ is infinitesimal. In R4, $1-H$ is infinitesimal
and $0<\overline T<1$. In particular
$H(1-H)\in\mathrm I_c$ in both rows.
The native inverse is
\[
r=\frac{j-p}{j+px},\quad h=\frac{1-HT}{1+HTx},\quad
a_s=\frac{j+px}{1+HTx},\quad
B_s=\frac{B(1+HTx)}{1+pB},\quad
t_s=\frac{T(1+pB)}{1+HTx}.                              \tag{EN4}
\]
The exact six weights are RSR6 at this substitution.
Those rational identities, including every contact first jet, were
checked against the native source without size restrictions. Their
former R1 size hypothesis is not imported.

Put $\sigma=1+(1-H)T$, $K=j+1-p$, $E=1+B(HT+p-1)$, and
use the explicit polynomials $\mathcal P,\mathcal L,\mathcal N,
\mathcal F,\mathcal K$ in RSR5, RSR15, RSR18. Then
\[
\frac{\mathcal P}{j+1}\sim1,\qquad
\frac{\mathcal L}{j+1}\sim\sigma,\qquad
\frac{\mathcal F}{j+1}\sim\sigma.                       \tag{EN5}
\]
Here $T,\sigma,(1-HT)^{-1}$ are bounded positive units.
To prove these uniform statements, divide the printed finite
polynomials by $j+1$. The corrections in $\mathcal P$ contain
$B+x$ or $TB$; those in $\mathcal L$ after
$(j+1)\sigma$ contain $p/(j+1)$ or $B+x$.
In $\mathcal F$ the only additional quotient is
$p(1-p)/(j(j+1))$, infinitesimal by (EN3).
Every other term contains $B$ or $Bj/(j+1)$.
Thus no restriction on $j$ or $jT$ is needed.

The exact physicality numerator has the grouping
\[
\mathcal N=p\{1+B(3j+1-p)+BT(j+1)\}
-Bj(j+1)\{1+(1-H)T\}+\mathcal R,\quad
|\mathcal R|\le C_0BT\,Bj(j+1).                         \tag{EN6}
\]
The bound follows termwise from RSR5 using $p<\min(1,j)$ and
bounded $T$. Since $n=\mathcal N/[2(j-p)\mathcal P]\ge0$,
\[
\begin{aligned}
&\eta_*=\frac{Bj(j+1)}p=O(1),\qquad
0\le n=O\!\left(\frac{p}{j(j+1)}\right)\in\mathrm I_c,\\
&\eta_*\in\mathrm I_c\Longrightarrow
n\sim\frac{p}{2(j-p)(j+1)}>0 .
\end{aligned}                                                    \tag{EN7}
\]
This includes $\mathcal N=0$ literally.

Write $\mu=m(1-p)$, $h_0=\mu+1/2$. The full source formulas give
\[
\begin{aligned}
&\mu=O(p),\quad h_0\asymp1,\quad
\alpha/[B^2(j+1)]\sim\sigma/2,\quad \alpha/B\in\mathrm I_c,\\
&A\sim\frac{T}{2(1-HT)},\qquad
CH\sim\frac{\sigma}{2T},\qquad
D(1-H)\sim\frac1{2T}.
\end{aligned}                                                    \tag{EN8}
\]
Consequently $A$ and the finite one of $C,D$ have positive real
lower bounds. The other coefficient is unbounded. The exact
normalized source limits are
\[
\begin{array}{c|c|c|c}
 & C&A&D\\ \hline
R3&+\infty&T/2&1/(2T)\\
R4&1/(2T)&T/[2(1-T)]&+\infty .
\end{array}                                                    \tag{EN9}
\]
Only real residues of the bounded entries are meant in this table.
The coefficient $m$ may be infinitesimal, a unit, or unbounded.

\subsection{The scalar allocation in the complete two-normal frame}

Use the entire parent RFN4 and the exact full-state change RFN5--RFN9,
with $w=(W+uY)/(1+u^2)=1$ and arbitrary $z=(Z-uX)/(1+u^2)>0$.
The identities are rational identities of (EN4), so they retain
both complete normals and every state term in the present source.
The tests in this subsection satisfy $u,v>0$ infinitesimal.
For clarity, the inverse is
\[
X=N_0+pv-uz,\quad Z=z+uN_0+puv,\quad
W=1+u(M_0+Hvz),\quad Y=u-M_0-Hvz.                       \tag{EN10}
\]
Retain the actual positive packet $mpD_u^2N_0^2$, and divide the
remaining entire parent by the pure testing factor $D_u=1+u^2$.
Its exact expansion is RFN7--RFN9. In particular the hard entries are
\[
\begin{aligned}
&H_N=h_0D_u+\tfrac12v^2+\alpha u^2+nD_v,\\
&H_M=D_u(CD_v+D+\tfrac12)+AD_v+(\alpha+\tfrac12)v^2,\qquad
H_{NM}=(\alpha+\tfrac12)uv .
\end{aligned}                                                    \tag{EN11}
\]

Keep the separated rows
$U=u/\sqrt{x}$, $V=v/\sqrt{y}$, $Z_*=z/\sqrt{q_z}$,
$Q_*=U-VZ_*$, $B_*=V-UZ_*$, $R_*=Z_*-UV$, and
\[
\Psi=U^2+V^2+Z_*^2+U^2V^2Z_*^2-4UVZ_*
=(U-V)^2+Z_*^2(1-UV)^2+2UV(1-Z_*)^2.                 \tag{EN12}
\]
The last radical is separated, for example
$\sqrt{UV}=(xy)^{-1/4}\sqrt{uv}$.
The exact scalar identity is RFN12 with the same actual circuit
masses. From (EN5), the closed formulas RSR16--RSR20 now give
\[
\begin{aligned}
&f_c\sim Tx^2\sigma/2,\quad
q_c\sim\frac{Tx}{2(1-HT)},\quad
r_c\sim TjB^2(1-p)\sigma/2,\\
&g_c/f_c=Hy\,\mathcal L/\mathcal F\in\mathrm I_c,\quad
k_c<0,\quad |k_c|/(q_cy)=O(p+B).
\end{aligned}                                                    \tag{EN13}
\]
In particular all three displayed positive masses are positive.
The sign of $k_c$ follows from the same exact positive-surplus
identity RSR19: its error bound uses $B(j+1)\in\mathrm I_c$,
not $T$ infinitesimal.

Before spending any current mass, perform the exact source-only
completion
\[
\begin{aligned}
h_0N_0^2+n(N_0+jb)^2
&=(n+h_0)(N_0+c_bb)^2+M_eb^2,\\
b&=v-uz/j,\quad c_b=\frac{nj}{n+h_0},\quad
M_e=\frac{nh_0j^2}{n+h_0}.
\end{aligned}                                                    \tag{EN14}
\]
Keep $nv^2(N_0+pb)^2$ and the positive extra hard terms RFN15
as whole squares. This reserves exactly the correlated $n$ packets
of the complete parent. If $n=0$, set $c_b=M_e=0$ and omit that
row. For $n>0$, $c_b^2/M_e=n/[h_0(n+h_0)]$ is bounded.

The signed scalar allocation differs from R1 in one source split.
The $e_0$ cost is unchanged: its negative part is at most $g_c$,
and RFN18b buys it at cost
$H(\mathcal L/\mathcal F)(y+v^2)\in\mathrm I$.
Also $d_c>0$. Use $Z_*(V^2-1)=-VQ_*-R_*$ for $k_c$.
The $VQ_*$ cost is $O((p+B)v^2)$, infinitesimal.
For its opposite-row part the exhaustive cases are:
\[
\begin{array}{c|l}
p\in\mathrm I_c&
 |k_c|/r_c=
 T\mathcal K/[(1-p)(1+Hy)\mathcal L]=O(p+B)\in\mathrm I_c\\
p\in\mathrm U_c,\ \eta_*\in\mathrm U_c&
 j\text{ unbounded},\ j^2B\in\mathrm U_c,\ 
 |k_c|/f_c\in\mathrm I_c,\ |k_c|/(f_cx)=O(1)\\
p\in\mathrm U_c,\ \eta_*\in\mathrm I_c&
 j\text{ unbounded},\ M_e\in\mathrm U_c,\ 
 |k_c|/f_c\in\mathrm I_c,\\
& |k_c|/(q_cy)=O(1),\ |k_c|/(M_eyx)=O(1).
\end{array}                                                    \tag{EN15}
\]
In the second case $j$ is unbounded because $p/j$ is infinitesimal;
$j^2B\asymp p\eta_*$. In the third,
$nj^2\sim p/2$ and $n$ is infinitesimal, so $M_e\sim p/2$.
The closed formulas show
$|k_c|/f_c=O((p+B)/j)$ in the last two cases and
$|k_c|/(M_eyx)=O((p+B)(1+1/n)/j^2)$ in the third.
These bounds are uniform through $p\to1$.

Case two uses RFN18a:
$R_*^2\preceq_{\rm sos}(1+U^2)\Psi$.
Case three uses RFN18c:
$R_*^2\preceq_{\rm sos}2(\Psi+V^2Q_*^2+U^2B_*^2)$.
The non-infinitesimal source ratios in (EN15) acquire the testing
factors $u^2$ or $v^2$. Thus all costs are infinitesimal.
Literal $n=0$ lies in the first or second case; it is not discarded.
After this finite allocation, fixed positive fractions of
\[
q_cQ_*^2,\quad r_cR_*^2,\quad f_c\Psi,\quad M_eb^2       \tag{EN16}
\]
remain, with the last row absent exactly at $n=0$.

\subsection{Both complete normal arrays, with the new endpoint weights}

The exact derivatives RFN21--RFN22, and their exact source cancellations
$S_N=S_0=S_1=S_2=0$, hold unchanged. Let
$q=vz-Tu$, $R=z-uv/B$, $R_u=v(u^2-x)$. In particular
\[
\begin{aligned}
L_N^{\rm rem}
&=-j\alpha b+\tfrac12pv^2b+K_RR_u+(K_c/j)uvq,\\
L_M&=C_qq+\mu uvb+v(\alpha-BTCHu^2)R,\\
K_c&=p/2+j\{H(\alpha+1/2)-1/2\}+n(p-j),\\
K_R&=-\mu-(\alpha+1/2)(1-p)+(T/j)K_c,\\
K_d&=CH-(D+1/2)(1-H)=-\mu/j-BTCH,\\
C_q&=K_d+AH+1/2+
v^2\{CH(1+u^2)+AH+(\alpha+1/2)H\}.
\end{aligned}                                                    \tag{EN17}
\]
Every displayed coefficient other than a retained row is bounded:
$p/j$ and $n$ are infinitesimal, $\mu/j=O(p/(j+1))$ is bounded,
and $A,CH,D(1-H),T$ are bounded by (EN8).
In particular $K_R$, $K_c/j$, and $C_q$ are bounded.

Here is an explicit audit of all costs, to avoid importing the former
$T\in\mathrm I_c$ hypothesis of RFN.
The exact radial Gram RFN19--RFN20 purchases $R_u^2$ at infinitesimal
cost from (EN16). Its three cases are still $1-p$ a unit;
$p\sim1,\eta_*$ a unit; and $p\sim1,\eta_*$ infinitesimal.
Their costs are respectively
\[
\begin{array}{c|l}
1-p\in\mathrm U_c&O(B),\ O(Bu^2),\ O(u^2/(1-p))\\
p\sim1,\eta_*\in\mathrm U_c&
 O(B+u^2/j+u^4/(j^2B)),\ O(Bu^2)\\
p\sim1,\eta_*\in\mathrm I_c&
 O(u^2/j),\ O(u^4/j),\
 [8u^4+2(u^2-x)^2]/M_e .
\end{array}                                                    \tag{EN18}
\]
They require only $T$ bounded above and below, because $T$ cancels
between $y=TB$ and the masses (EN13).
The small-current budgets NSR3--NSR5 likewise apply unchanged:
both $pv^2b$ and $\sqrt T\,\mu uvb$ are free.
For $\eta_*$ infinitesimal their coefficients are
$O(v^4)$ and $O(Tu^2v^2)$, using $p^2/M_e=O(1)$.
For $\eta_*$ a unit the coefficients are
$O(v^2(y+v^2))$ and $O(Tu^2(y+v^2))$, using
$p^2y/f_c=O(1)$. Each is infinitesimal.

The row $q/\sqrt T$ is strict by (EN13).
Its $N$ coupling has coefficient $(K_c/j)uv\sqrt T$ infinitesimal.
The identity $j\alpha b=-(\alpha/B)R_u-\alpha uR$ handles
the remaining $N$ term. For $1-p$ a unit the latter square costs
$O(B^2(j+1)^2u^2/(1-p))$ against $r_cR_*^2$.
For $p\sim1$, RFN18a bounds the entire $j\alpha b$ square by
$O(B^3(j+1)^2+B^2(j+1)^2u^2/j)$ of the circuit.
These are infinitesimal. A sufficiently large fixed real radial
pivot pays the bounded $K_R$ coupling, including the bounded shift
$c_b/\sqrt{M_e}$ from (EN14).

Normalize the complete second normal by
$M_0=\sqrt{TH(1-H)}\,\widehat M$. The full hard diagonal (EN11)
then has residue
\[
\overline{TH(1-H)H_M}=
\begin{cases}
(1+\overline T)/2,&R3,\\
1/2,&R4 .
\end{cases}                                                    \tag{EN19}
\]
Every term of the complete diagonal is bounded: the only divergent
factors occur in $CH$ and $D(1-H)$.
The normalized $q$ coupling is
$T\sqrt{H(1-H)}C_q$, now infinitesimal because
$H(1-H)$ is infinitesimal. This is the replacement for the
$T$-infinitesimal argument in R1.
The normalized current coupling is a bounded multiple of the
already purchased $\sqrt T\,\mu uvb$.

For the entire remaining opposite row put $J=B(j+1)+u^2$ and
$a_R=Tq_zv^2(\alpha-BTCHu^2)^2$.
Since $\alpha-BTCHu^2=O(BJ)$, its costs are
\[
\begin{aligned}
a_R/f_c&=O((T/j)v^2J^2),&
a_R/(q_cy)&=O(v^2J^2),\\
a_R/r_c&=O(Tv^2J^2/(1-p)),&
a_R/(M_eyx)&=O(T(1+1/n)v^2J^2/j^2).
\end{aligned}                                                    \tag{EN20}
\]
The direct $r_c$ cost applies for $1-p$ a unit.
For $p\sim1,\eta_*$ a unit, RFN18a applies using $j^2B$ a unit.
For $p\sim1,\eta_*$ infinitesimal, RFN18c applies using
$(1+1/n)/j^2=O(1)$. Thus the entire square, including every
higher testing degree, is free. No state ratio occurs in a cost.

Finally the complete mixed hard term is
$2(\alpha+1/2)uv\sqrt{TH(1-H)}
(\widehat N/\sqrt{n+h_0}-c_bb)\widehat M$.
Its bounded coefficients have zero residue.
The strict bounded-Gram argument therefore applies exactly as in
the final paragraph of RFN: choose a large fixed radial pivot, buy
the finitely many free rows, and retain positive fractions of
(EN16). The two hard pivots are $1$ and the positive number
(EN19). Restore all whole packets in (EN14), RFN15, and
$mpD_u^2N_0^2$, and the pure testing factor $D_u$.
This proves the full local parent for all infinitesimal $u,v>0$
and arbitrary $z>0$, with both unrestricted state normals.
The exact SJ state/sign reduction supplies every state and the
literal axes in this local chart.

\subsection{Every other testing angle, with weak \(m\) retained}

This outer join uses only (EN9), positivity of the six native
weights, and exact full-parent identities. Set
$\widetilde m=m/(1+m)$, a bounded positive source coefficient below
$m$, and retain the difference as the whole original $m$ square.
Keep the entire $n$ and $\alpha$ events throughout.
For the single unbounded coefficient $C$ or $D$, retain its excess
over a sufficiently large positive real cap. Other selected
coefficients have the actual bounded positive residues (EN9).

\paragraph{Finite Bob with nonzero residue and bounded Alice.}
Let $v$ have nonzero finite residue, $u$ be bounded, and use the
five full-state rows
\[
P=uW,\quad Y,\quad X,\quad Z,\quad S=\sqrt{\widetilde m}\,W.
                                                               \tag{EN21}
\]
The selected $m$ event is the entire square
$[vS-\sqrt{\widetilde m}(X-uvY+uZ)]^2$.
All other terms use $P,Y,X,Z$ with bounded coefficients.
In R3 the strong row is $P-Y$; on its kernel the selected remaining
form on $(Y,X,Z)$ is
\[
\frac T2D_uD_vY^2+\frac{1+T}{2T}v^2(uX-Z)^2+K_0 .
\]
In R4 the strong row is $P+uvX-Y-vZ$; on its kernel it is
\[
\frac{D_vv^2}{2T}(Z-uX)^2+
\frac{T}{2(1-T)}D_uD_vY^2+K_0 ,
\quad
K_0=\frac{D_u(vY-Z)^2+D_v(X+uZ)^2-D_uD_vZ^2}{2}.
\]
The leading principal minors in the order $(Y,X,Z)$ are
\[
\begin{array}{c|c|c}
 &R3&R4\\ \hline
1&
\dfrac{D_u(TD_v+v^2)}2&
\dfrac{D_u(T+v^2)}{2(1-T)}\\[4pt]
2&
\dfrac{D_u(TD_v+v^2)(Tu^2v^2+Tv^2+T+u^2v^2)}{4T}&
\dfrac{D_uD_v(T+v^2)(T+u^2v^2)}{4T(1-T)}\\[4pt]
3&
\dfrac{v^2(1+T)D_u^2D_v(Tu^2+v^2)}{8T}&
\dfrac{v^2D_u^2D_v^2(Tu^2+v^2)}{8T(1-T)} .
\end{array}                                                    \tag{EN22}
\]
All are strictly positive at the real residues, even $u=0$.
The $S$ direction has the positive coefficient $v^2$ in the
selected $m$ event. A sufficiently large real cap on the strong
row makes the entire five-row residue positive definite.
This is the real finite-cap Schur lemma: if a bounded real
quadratic form is positive on the kernel of a linear map, adding
a sufficiently large real multiple of the squared map makes it
positive definite. The unbounded actual coefficient permits
every such cap. This argument includes $m$ infinitesimal.

\paragraph{Infinitesimal Bob and Alice a finite positive unit.}
Use $R=uW-Y$, $B_t=X+uZ$, $V_Z=vZ$, and $Y$; the testing inverse
$1/u$ is bounded. The full selected parent is a bounded quadratic
on these four rows because
$v(W+uY)=(v/u)(R+D_uY)$ and
$J=R+uvB_t-D_uV_Z$.
Its residue splits off $(\widetilde m+1/2)B_t^2$ and the form
\[
CR^2+AD_uY^2+(D+1/2)(R-D_uV_Z)^2
-D_uYV_Z-\tfrac12D_uV_Z^2.                            \tag{EN23}
\]
In R3 restrict to the strong-row kernel $R=0$.
After removing a positive factor $D_u$, the matrix on $(Y,V_Z)$ is
$\left(\begin{smallmatrix}T/2&-1/2\\-1/2&
(1+T)D_u/(2T)-1/2\end{smallmatrix}\right)$;
its determinant is $(1+T)u^2/4>0$.
In R4 the strong kernel is $R=D_uV_Z$.
On $(R,Y)$ its first diagonal is $1/(2T)-1/(2D_u)>0$,
and its determinant is $u^2/[4(1-T)]>0$.
A large real cap again makes (EN23) strict.

\paragraph{Infinite Bob and bounded Alice.}
Put $w=1/v$, divide the homogeneous parent by $v^2$, and use
the five rows (EN21). Select any fixed positive real lower caps
$c_0,a_0,d_0$ below $C,A,D$, and the actual $\widetilde m$ event.
The full quadratic is bounded; its residue is
\[
[S+\sqrt{\widetilde m}\,uY]^2+
c_0(P-Y)^2+(a_0+\tfrac12)D_uY^2+
d_0(uX-Z)^2+\tfrac12D_uX^2.                           \tag{EN24}
\]
This is positive definite on the five rows, for every bounded
residue of $u$. It includes $u=0$ and $m$ infinitesimal.

\paragraph{Infinite Alice.}
Let $s=1/u$ and keep the entire $m,n,\alpha$ events.
For finite $v$ of nonzero residue, after division by $u^2$
the selected residue on $(W,Y,X,Z)$ is
\[
c_0D_vW^2+a_0D_vY^2+
(d_0+\tfrac12)(W+vX)^2+\tfrac12(vY-Z)^2,
\]
which is positive definite. For $v$ infinitesimal use instead
$(N,W,Y,V)=(sX+Z,W,Y,vX)$.
The complete signed core after division by $u^2$ is
\[
\frac12\{W+(1+s^2)V-sY-svN\}^2
+\frac12(1+s^2)v^2Y^2
-(1+s^2)Y(vN-sV)-\frac12(1+s^2)(vN-sV)^2
+\frac12D_vN^2 .
\]
Adding the selected $C,A,D$ packets gives strict residue
$c_0W^2+a_0Y^2+(d_0+1/2)(W+V)^2+N^2/2$.
Every full coefficient is bounded; no $s/v$ is used.
If also Bob is infinite, divide by $u^2v^2$ and use
$(W,X,Y,S_Z,T_Z)=(W,X,Y,sZ,wZ)$.
The complete selected residue is
\[
c_0W^2+a_0Y^2+d_0(X-S_Z)^2+\tfrac12X^2
+\tfrac12(Y-T_Z)^2,                                   \tag{EN25}
\]
again positive definite. The full quadratic follows directly from
$J/(uv)=wW+X-swY-S_Z$ and the exact signed core;
its coefficients are bounded polynomials in $s,w$.

These charts exhaust all nonnegative projective angles:
both infinitesimal are the complete local theorem, and every other
pair is one of (EN21)--(EN25) or (EN23).
All strict-matrix coefficients use only bounded pure source or
pure testing factors. Restore every selected excess, the untouched
$m,n,\alpha$ packets, and the homogeneous testing squares.
The exact SJ sign gauges, literal projective endpoints and
real/imaginary state split finish the asserted all-tests theorem.

\paragraph{Audit scope and dependencies.}
The source dictionary and exact scalar/normal identities are the
freshly reconstructed RSR, NSR, RFN sources; only their explicitly
rechecked algebraic identities and finite SOS budgets are used.
Equations (EN3)--(EN20) prove all changed source and coefficient
bounds afresh. The entire outer join (EN21)--(EN25) is new.
The focused checker corroborates source coefficient limits,
full-parent expressions, complete hard entries, and every printed
positive kernel matrix. Analytical uniform estimates and the
exhaustive join are the proof above, not an inference from its receipt.

